\chapter{Proposed Work: Design and Implementation of the Memory Layer}

\section{Overview of the Methodology}
Phase 1 proposed a memory layer that (i) converts user interactions into embeddings, (ii) stores them in a
vector database, (iii) retrieves the top-$k$ relevant memories at query time and (iv) injects them into the
prompt. Phase 2 implements this design and extends it with the components listed as future work in Phase 1.
Figure~\ref{fig:arch} shows the resulting architecture. Conversation turns enter on the left through the
\emph{write path}; questions enter through the \emph{read path}, which is a multi-stage retrieval pipeline
ending in a grounding gate.

\begin{figure}[htbp]
\centering
\begin{tikzpicture}[x=1cm,y=1cm,
  box/.style={draw=black!60,fill=boxfill,rounded corners=2pt,align=center,font=\footnotesize,
              minimum height=9mm,inner sep=2pt,text width=27mm},
  store/.style={draw=black!60,fill=white,cylinder,shape border rotate=90,aspect=0.25,align=center,
                font=\footnotesize,minimum height=10mm,text width=22mm,inner sep=2pt},
  lbl/.style={font=\scriptsize\itshape,text=black!60},
  arr/.style={-{Stealth[length=2mm]},thick,draw=black!70}
]
% ---- write path (top row) and the three stores
\node[box] (turn) at (0,0) {Conversation turn\\(user / assistant)};
\node[box] (norm) at (3.6,0) {Normalisation\\(third person)};
\node[box] (imp)  at (7.2,0) {Importance scoring\\$+$ pruning};
\node[box] (ent)  at (10.8,0) {Entity\\extraction};
\node[store] (vec) at (3.6,-2.1) {FAISS\\vector index};
\node[store] (bm)  at (7.2,-2.1) {BM25\\keyword index};
\node[store] (gr)  at (10.8,-2.1) {Entity\\graph};
\draw[arr] (turn) -- (norm); \draw[arr] (norm) -- (imp); \draw[arr] (imp) -- (ent);
\draw[arr] (norm) -- node[lbl,left]{embed} (vec);
\draw[arr] (imp)  -- node[lbl,left]{tokens} (bm);
\draw[arr] (ent)  -- node[lbl,left]{entities} (gr);
\node[lbl,anchor=south west] at (-1.45,0.55) {Write path};
% ---- read path, first row (left to right)
\node[box] (q)    at (0,-4.4) {Question};
\node[box] (fuse) at (3.6,-4.4) {Weighted RRF\\(dense $+$ sparse,\\query expansion)};
\node[box] (gx)   at (7.2,-4.4) {Graph expansion\\(1 hop)};
\node[box] (rr)   at (10.8,-4.4) {Cross-encoder\\re-ranking};
% ---- read path, second row (right to left)
\node[box] (sc)   at (10.8,-6.5) {Memory scoring $+$\\conflict-aware recency};
\node[box] (ctx)  at (7.2,-6.5) {Top-$k$ notes\\(token budget)};
\node[box] (llm)  at (3.6,-6.5) {LLM\\(Qwen2.5-1.5B)};
\node[box] (gate) at (0,-6.5) {Grounding\\gate};
\node[box] (ans)  at (0,-8.4) {Answer or\\``I don't know''};
\draw[arr] (q) -- (fuse); \draw[arr] (fuse) -- (gx); \draw[arr] (gx) -- (rr);
\draw[arr] (rr) -- (sc); \draw[arr] (sc) -- (ctx); \draw[arr] (ctx) -- (llm);
\draw[arr] (llm) -- (gate); \draw[arr] (gate) -- (ans);
\draw[arr,dashed] (vec) -- (fuse); \draw[arr,dashed] (bm) -- (fuse); \draw[arr,dashed] (gr) -- (gx);
\node[lbl,anchor=south west] at (-1.45,-3.85) {Read path};
\end{tikzpicture}
\caption{Architecture of the implemented memory layer. Dashed arrows are index look-ups.}
\label{fig:arch}
\end{figure}

\section{Data Acquisition and Benchmark Design}
\label{sec:data}
Public long-term memory benchmarks (LoCoMo, LongMemEval) are designed for large API models and do not give
ground-truth \emph{memory identifiers}, which we need to measure retrieval recall. We therefore generated a
controlled benchmark in the same spirit (\texttt{benchmarks/make\_dataset.py}). Each simulated user talks
to an assistant over 12 sessions of 12 user--assistant pairs (288 turns, about 3.6k tokens of history).
Into the small talk we insert:
\begin{itemize}[nosep]
  \item \textbf{10 single-hop personal facts} drawn from 18 attribute types (city, job, pet, allergy, doctor,
        laptop, \dots), each stated once with one of three paraphrases.
  \item \textbf{2 updated facts} (e.g.\ a move or a job change) stated in a later session, which create a
        \emph{temporal-current} question (``Where does the user live now?'') and a \emph{temporal-before}
        question (``Where did the user live before moving?'').
  \item \textbf{2 multi-hop chains} in which a related person is introduced in one session and a fact about
        that person appears in another (``My sister Priya visited'' $\rightarrow$ ``Priya adopted a cat named
        Milo'' $\rightarrow$ ``What did the user's sister name her cat?'').
  \item \textbf{8 distractor facts} about other people that reuse the attribute vocabulary (``My neighbour's
        dog is called Luna'') to punish lexical matching.
  \item \textbf{3 unanswerable questions} about attributes that were never mentioned; the correct answer is
        an abstention.
\end{itemize}
Questions are paraphrased relative to the stored statements, so a question rarely shares its exact wording
with its evidence. The test split contains 20 users and 330 questions (170 single-hop, 30 temporal-current,
30 temporal-before, 40 multi-hop, 60 unanswerable). A separate 6-user development split with a different
random seed was used for all hyper-parameter tuning. A second generator builds the NIAH-style histories of
1k--16k tokens used in Section~\ref{sec:niah}.

\section{Write Path}
\subsection{Memory Normalisation}
\label{sec:normalisation}
Questions are asked about ``the user'' while memories are written in the first person. Both the embedding
model and the cross-encoder treat ``I live in Pune'' and ``Where does the user live?'' as only weakly
related. Each turn is therefore rewritten into the third person at write time with a small set of pronoun
rules (``I'' $\rightarrow$ ``the user'', ``my'' $\rightarrow$ ``the user's'', assistant ``you'' $\rightarrow$
``the user''). On a probe set this raised the cosine similarity of question--memory pairs from about 0.3 to
0.6 and moved the cross-encoder logit of the correct memory from $-9$ (confidently irrelevant) to $+4$.

\subsection{Importance Scoring and Pruning}
Every turn receives an importance score in $[0,1]$ from cheap surface cues: user turns score higher than
assistant turns; statements that mention the speaker, contain a capitalised entity or a number, use
personal-fact vocabulary (\emph{live, work, named, allergic, birthday, deadline, \dots}) or are longer than
five words gain points, and questions lose points. A configurable threshold discards low-importance turns at
write time; Section~\ref{sec:pruning} shows that about 84\% of the turns can be discarded without any loss of
retrieval recall.

\subsection{Indexes and Entity Graph}
Each stored memory is embedded with \texttt{all-MiniLM-L6-v2} \cite{sbert} (384 dimensions, L2-normalised)
and inserted into a flat inner-product FAISS index \cite{faiss}; its stemmed, stop-word-filtered tokens are
added to a BM25 index \cite{bm25}; and its entities (capitalised names and product names, dates, numbers,
extracted with regular expressions and a sentence-initial filter) become nodes in a bipartite memory--entity
graph implemented with NetworkX. Entities linked to more than 25 memories are treated as hubs and ignored
during expansion.

\section{Read Path}
\subsection{Hybrid First Stage and Fusion}
\label{sec:fusion}
The question is embedded and the 20 nearest memories are fetched from FAISS. In parallel, the question is
expanded with a rule-based synonym table (\emph{city} $\rightarrow$ \emph{live, based, moved, home, \dots})
and the 20 best BM25 matches are fetched; BM25 hits scoring below 20\% of the best hit are discarded, which
removes near-zero ``matches'' caused by frequent words. The two lists are merged by weighted reciprocal
rank fusion \cite{rrf}
\[
  s_{\text{fused}}(m) = \sum_{c \in \{\text{dense},\,\text{sparse}\}} \frac{w_c}{k + \text{rank}_c(m)},
  \qquad k = 60,\; w_{\text{dense}} = 1,\; w_{\text{sparse}} = 0.2 .
\]
The sparse weight was chosen on the development split: with equal weights the weak lexical channel reduced
recall below that of dense retrieval alone.

\subsection{Entity-Graph Expansion}
The top-3 fused candidates act as seeds. Every memory that shares a non-hub entity with a seed inherits
$0.85\times$ the seed's fused score (one-step relevance propagation); at most four such memories are added
to the candidate pool. This is how the second hop of a chain question enters the pool: ``Dev started working
at Zoho'' never mentions a brother, but it shares the entity \emph{Dev} with ``The user's brother Dev lives
with the user'', which the first stage finds easily.

\subsection{Cross-Encoder Re-ranking}
Up to 40 candidates are scored jointly with the question by the MS~MARCO MiniLM-L6 cross-encoder
\cite{nogueira2019}. Because the cross-encoder was found to be unreliable on its own for these short
conversational notes, its sigmoid output is blended with the normalised fused score,
$\text{rel}(m) = 0.5\,\sigma(\text{CE}(m)) + 0.5\, s_{\text{fused}}(m)/\max s_{\text{fused}}$.

\subsection{Memory Scoring}
\label{sec:scoring}
Following the memory-stream idea \cite{park2023}, each candidate is scored with relevance, recency,
importance, access frequency and a graph-coherence bonus (number of other candidates it shares an entity
with). Unlike the additive formulation, relevance acts as a gate:
\[
  \text{score}(m) = \text{rel}(m)\,\bigl(1 + w_{\text{rec}}\,2^{-\Delta s/h} + w_{\text{imp}}\,\text{imp}(m)
  + w_{\text{acc}}\,\text{acc}(m) + w_{\text{graph}}\,\text{coh}(m)\bigr),
\]
with $\Delta s$ the age in sessions and half-life $h = 6$. On the development split a non-zero recency
weight lowered recall for multi-hop and ``before'' questions, so recency is applied through the conflict
rule below instead ($w_{\text{rec}} = 0$, $w_{\text{imp}} = 0.10$, $w_{\text{acc}} = 0.02$,
$w_{\text{graph}} = 0.10$). Finally, a memory that is itself a question or request to the assistant
(``What is a good way to learn guitar chords?'') is multiplied by $0.7$: such turns are rarely evidence for
a fact about the user, yet they match question words such as \emph{learning} very well. This single rule
raised single-hop Recall@5 from 92\% to 99\% on the test split without affecting the other question types.

\subsection{Conflict-Aware Recency (Memory Updates)}
Two candidates are treated as versions of the same fact when both are user statements about the user (no
third-party nouns such as \emph{colleague} or \emph{neighbour}), their cosine similarity is at least 0.45 and
the newer one contains an update cue (\emph{moved, switched, upgraded, replaced, now, new, \dots}). The older
version keeps 80\% of its score and is tagged ``(older note)'' when injected, so the LLM can still answer
``before'' questions while being steered to the current fact.

\subsection{Context Injection and Grounding Gate}
The top-5 memories are rendered chronologically with their session label inside a 512-token budget, and the
LLM is instructed to answer only from the notes, to prefer the most recent note when notes conflict, to
ignore notes about other people, and to reply ``I don't know'' otherwise. After generation the
\emph{grounding gate} verifies the answer: the note whose content covers the answer's content words is
located, and the answer is accepted only if that note's relevance score is at least 0.35. An answer copied
from an irrelevant note (the neighbour's dog) or not found in any note is replaced by an abstention. This
gate is the Phase-2 counterpart of the hallucination detectors evaluated in Phase 1.

\section{Baselines and Evaluation Protocol}
All strategies use the same model, Qwen2.5-1.5B-Instruct \cite{qwen25} in fp16 on a 4~GB laptop GPU, with
greedy decoding, the same system prompt and the same answer-format instructions:
\begin{itemize}[nosep]
  \item \textbf{Baseline LLM}: question only, no memory.
  \item \textbf{Full context}: the complete transcript of all sessions in the prompt (the context-window
        expansion approach; prompts up to 16k tokens use chunked prefill so that they fit in GPU memory).
  \item \textbf{Sliding window}: the last 40 turns.
  \item \textbf{Vector RAG}: the Phase-1 design, embedding similarity top-5 only.
  \item \textbf{Memory layer}: the full pipeline above, reported with and without the grounding gate.
\end{itemize}
An answer is correct when the normalised gold string is contained in it or token-F1 $\ge 0.6$; an
unanswerable question is correct only when the model abstains. A \emph{hallucination} is a non-abstaining
answer that is wrong, or any non-abstaining answer to an unanswerable question. For retrieval methods we also
report Recall@$k$ (all gold memories within the top-$k$) and MRR. Latency and prompt tokens are measured per
question. The implementation is about 2,400 lines of Python: the package \texttt{memory\_layer} (eleven modules),
the \texttt{benchmarks} directory (dataset generator, three experiment runners, tuning and plotting) and a
functional test suite; all experiments are reproducible from the committed scripts.
